\section{Conclusion}
\label{sec:conclusion}

This paper applied Edmonds-Karp max-flow to the ROUTE national directed graph of 227 corridors and approximately 48,000 directed edges to identify the binding capacity constraints in the US interstate freight network and to quantify the throughput effects of incident scenarios and missing link completions.

\paragraph{Three binding bottlenecks.} The national minimum cut analysis identifies three arcs that bind multiple O-D cluster pairs simultaneously: the I-95 Baltimore-Washington segment (210,000 vpd on 6-lane facility; V/C 2.1+ at peak; \$1.4B/yr congestion cost), the I-80 Donner Pass crossing (91,200 vpd maximum; geographic chokepoint with no interstate alternate; 50 closures per year), and the Dallas I-35/I-45 interchange complex (V/C 1.9+; binding for all Texas Gulf-to-Midwest and Gulf-to-Northeast flows). These three arcs account for the binding minimum cut in four of six major cross-regional O-D pairs.

\paragraph{Donner closure cascades.} The Donner Pass incident simulation confirms that a multi-day closure is not an isolated corridor event but a national throughput shock. NE-to-Pacific max-flow drops 23 percent (from 320,000 to 246,000 vpd); the redirected flow pushes I-40 from V/C 0.52 to 0.84, within 12,000 vpd of saturation. The simultaneous Donner-plus-I-35-Oklahoma compound scenario drives I-40 to V/C 1.11---a network failure condition in which the national T2 resilience layer is itself saturated. The annualized cost of Donner-related throughput suppression and detour is estimated at \$3.4 billion per year \citep{ATRI_costs2024}.

\paragraph{I-69 is the highest-impact missing link.} I-69 completion (approximately 400 missing miles, Texarkana to Indianapolis) increases Texas Gulf-to-Midwest max-flow by 18 percent by creating a direct path that bypasses the Dallas binding constraint. Dallas V/C drops from 1.9+ to approximately 1.4. The NPV is negative at 7 percent but approaches zero when freight growth is included and is positive under social discount rate assumptions. Of the three missing links analyzed, I-69 produces the largest single max-flow gain and addresses the most congested interchange node in the network.

\paragraph{I-70W is the highest-value resilience asset.} I-70W (Cove Fort, UT to Sacramento, CA; 700 miles) contributes only 4.7 percent to baseline NE-Pacific max-flow but reduces the Donner closure throughput loss from 23 percent to 9 percent. Its NPV is positive at 5 percent when resilience benefits are credited (\$6.7B net present value). I-70W is justified not as a capacity addition under normal operations but as insurance against the compound failure scenario described above.

\paragraph{National capacity target.} The current network maximum of 4.8 million commercial vehicles per day must reach 7.2 million by 2050, a 50 percent increase driven by the FHWA 1.8 percent per year freight growth projection \citep{FHWA_freight2023}. This target is achievable through a combination of managed lane additions on the 7 T1 Primary Arterial corridors with binding minimum cuts, I-69 completion, and I-70W construction. No single investment achieves the target; the combination is necessary because different O-D pairs have different binding arcs.

\paragraph{Limitations.} This analysis uses single-commodity max-flow, treating all truck traffic as equivalent. A multi-commodity formulation would differentiate by freight class, value, and time-sensitivity---allowing the model to distinguish, for example, between time-sensitive pharmaceutical shipments that cannot tolerate Donner-level delays and bulk commodity traffic that can use I-40 without shipper contract penalties. The single-commodity formulation understates the economic cost of constraints on high-value freight and overstates the resilience of rerouting options for time-sensitive freight.

The capacity estimates also rely on HPMS 2018 data \citep{HPMS2018}, now six years old. Several corridor segments, particularly in high-growth states (Texas, Florida, Tennessee), will have changed substantially. Updated HPMS data, when available, should be applied to refresh the bottleneck ranking.

\paragraph{Next: D.1 climate exposure.} The incident scenarios in Section~\ref{sec:incident} treat Donner closures and I-35 disruptions as stationary processes---constant frequency and duration across the analysis period. The D-track analysis begins with D.1, which maps climate exposure of the interstate network through 2050. If Sierra Nevada precipitation intensity increases with warming, Donner closure frequency increases; if Oklahoma flooding events increase in frequency, I-35 disruptions become more common. The compound failure scenario identified in this paper is not hypothetical in 2050 conditions; D.1 will quantify the probability that it becomes a recurring seasonal event rather than a rare coincidence.
